% =====================================================================
% A Multimodal AI Framework for Detecting Coordinated Propaganda
% in Bangla Online Media Content
% B.Sc. in CSE, United International University (UIU)
% Compile with:  xelatex report.tex   (run twice for cross-references)
% =====================================================================
\documentclass[12pt,a4paper]{report}

% ---------- Fonts & Engine Compatibility (XeLaTeX or pdfLaTeX) ----------
\usepackage{iftex}
\ifXeTeX
  \usepackage{fontspec}
  \setmainfont{TeX Gyre Termes}[Ligatures=TeX]
  \setsansfont{TeX Gyre Heros}[Ligatures=TeX]
  \setmonofont{TeX Gyre Cursor}[Scale=0.9]
  % Bangla text (used in UI mock-up)
  \newfontfamily\banglafont{FreeSerif}
  \newcommand{\bn}[1]{{\banglafont #1}}
\else
  \usepackage[T1]{fontenc}
  \usepackage[utf8]{inputenc}
  \usepackage{tgtermes}
  \usepackage{tgheros}
  \usepackage{tgcursor}
  \newcommand{\bn}[1]{[Bangla text: \textit{ekhane bangla social media post ba khobor likhun\dots}]}
\fi

% ---------- Layout ----------
\usepackage[a4paper,left=1.5in,right=1in,top=1in,bottom=1in,headheight=15pt]{geometry}
\usepackage{setspace}
\onehalfspacing
\usepackage{fancyhdr}
\usepackage{titlesec}
\usepackage{enumitem}
\setlist{itemsep=2pt,topsep=4pt}
\usepackage{caption}
\captionsetup{font=small,labelfont=bf,skip=6pt}
\usepackage{float}
\usepackage{adjustbox}   % put next to \usepackage{float}
\usepackage{pdflscape}

% ---------- Tables ----------
\usepackage{array}
\usepackage{booktabs}
\usepackage{longtable}
\usepackage{tabularx}
\usepackage{multirow}
\usepackage{makecell}
\usepackage{ragged2e}
\newcolumntype{L}[1]{>{\RaggedRight\arraybackslash}p{#1}}
\newcolumntype{C}[1]{>{\Centering\arraybackslash}p{#1}}

% ---------- Math, symbols, colour ----------
\usepackage{amsmath,amssymb}
\usepackage{xcolor}
\definecolor{uiublue}{RGB}{0,51,102}
\definecolor{boxfill}{RGB}{232,240,250}
\definecolor{boxdraw}{RGB}{0,51,102}
\definecolor{storefill}{RGB}{255,244,214}
\definecolor{procfill}{RGB}{225,242,232}
\definecolor{extfill}{RGB}{244,228,228}

% ---------- Diagrams ----------
\usepackage{tikz}
\usetikzlibrary{arrows.meta,positioning,shapes.geometric,shapes.misc,fit,backgrounds,calc}
\usepackage{pgfgantt}

% ---------- Code / mock-ups ----------
\usepackage{listings}
\lstset{basicstyle=\ttfamily\footnotesize,breaklines=true,frame=single,
  columns=fullflexible,keepspaces=true,xleftmargin=0pt,framexleftmargin=4pt}

% ---------- Bibliography & links ----------
\usepackage[numbers,sort&compress]{natbib}
\usepackage[hidelinks,bookmarks=true]{hyperref}
\usepackage{url}
\urlstyle{same}

% ---------- Headers ----------
\pagestyle{fancy}
\fancyhf{}
\fancyhead[L]{\small\itshape Bangla Propaganda Detection Framework}
\fancyhead[R]{\small\thepage}
\renewcommand{\headrulewidth}{0.4pt}
\fancypagestyle{plain}{\fancyhf{}\fancyfoot[C]{\thepage}\renewcommand{\headrulewidth}{0pt}}

% ---------- Chapter/section styling ----------
\titleformat{\chapter}[display]
  {\normalfont\huge\bfseries\color{uiublue}}
  {\chaptertitlename\ \thechapter}{14pt}{\Huge}
\titlespacing*{\chapter}{0pt}{-20pt}{24pt}
\titleformat{\section}{\normalfont\Large\bfseries\color{uiublue}}{\thesection}{0.8em}{}
\titleformat{\subsection}{\normalfont\large\bfseries}{\thesubsection}{0.8em}{}
\titleformat{\subsubsection}{\normalfont\normalsize\bfseries\itshape}{\thesubsubsection}{0.8em}{}

% ---------- Convenience macros ----------
\newcommand{\ck}{$\checkmark$}
\newcommand{\ttl}{A Multimodal AI Framework for Detecting Coordinated Propaganda in Bangla Online Media Content}

% ---------- TikZ styles shared by all figures ----------
\tikzset{
  >=Stealth,
  ext/.style   ={rectangle,draw=boxdraw,fill=extfill,rounded corners=2pt,
                 minimum height=12mm,text width=24mm,inner sep=3pt,align=center,font=\scriptsize},
  proc/.style  ={rectangle,draw=boxdraw,fill=procfill,rounded corners=3pt,
                 minimum height=12mm,text width=24mm,inner sep=3pt,align=center,font=\scriptsize},
  blk/.style   ={rectangle,draw=boxdraw,fill=boxfill,rounded corners=3pt,
                 minimum height=12mm,text width=24mm,inner sep=3pt,align=center,font=\scriptsize},
  store/.style ={rectangle,draw=boxdraw,fill=storefill,
                 minimum height=12mm,text width=24mm,inner sep=3pt,align=center,font=\scriptsize},
  flow/.style  ={->,thick,draw=boxdraw!80!black},
  dflow/.style ={->,thick,dashed,draw=boxdraw!80!black},
  line/.style  ={thick,draw=boxdraw!80!black},
  lab/.style   ={font=\scriptsize,inner sep=1.5pt,align=center,text=black},
  grp/.style   ={draw=boxdraw!60,dashed,thick,rounded corners=4pt},
  grptitle/.style={font=\scriptsize\bfseries,text=boxdraw,inner sep=2pt}
}

\hypersetup{pdftitle={\ttl},pdfauthor={Group 262-037, UIU}}

\begin{document}

% =====================================================================
%                              TITLE PAGE
% =====================================================================
\begin{titlepage}
\thispagestyle{empty}
\centering
{\Large\bfseries\color{uiublue} United International University\par}
\vspace{2pt}
{\large Department of Computer Science and Engineering\par}
\vspace{1.4cm}
{\LARGE\bfseries\color{uiublue}\ttl\par}
\vspace{1.4cm}
{\large A Final Year Design Project (FYDP) Report submitted in partial fulfillment of the
requirements for the degree of\par}
\vspace{6pt}
{\large\bfseries Bachelor of Science in Computer Science and Engineering\par}
\vspace{1.2cm}
{\large\bfseries Submitted by\par}
\vspace{6pt}
\begin{tabular}{l c}
Shekh Abdullah Al Mehedi & ID: 112231075\\[3pt]
Mahjabin Khan & ID: 112230177\\[3pt]
Omor Faruck Ullas & ID: 112310384\\[3pt]
Md.\ Rayan Rahman & ID: 112331071
\end{tabular}
\vspace{1.2cm}

{\large\bfseries Group: 262-037\par}
\vspace{1cm}
{\large\bfseries Supervisor\par}
\vspace{4pt}
{\large Dr.\ Jannatun Noor Mukta\par}
\vfill
{\large Dhaka, Bangladesh\par}
{\large September 2026\par}
\end{titlepage}

\pagenumbering{roman}

% =====================================================================
%                              ABSTRACT
% =====================================================================
\chapter*{Executive Summary and Abstract}
\addcontentsline{toc}{chapter}{Executive Summary and Abstract}

Bangladesh's digital media ecosystem has experienced exponential expansion,
currently hosting over 77.7~million internet users and 60.0~million social media
user identities. While digital platforms democratize information exchange, they
also serve as critical conduits for strategic, coordinated propaganda campaigns
aimed at shaping public opinion and destabilizing civic discourse. Traditional
automated propaganda detection architectures predominantly focus on high-resource
languages (e.g., English) and evaluate isolated, text-only posts. In contrast,
under-resourced languages like Bangla lack comprehensive multimodal annotated
corpora, fine-grained persuasive technique taxonomies, and coordinated campaign
detection systems.

This report presents the design and architectural methodology of a
\textbf{Multimodal AI Framework for Detecting Coordinated Propaganda in Bangla
Online Media Content}. The proposed framework unites:
\begin{enumerate}
  \item A hybrid human-in-the-loop annotation pipeline utilizing Large Language
        Models (LLMs) for pre-annotation and domain experts for ground-truth
        verification.
  \item An extensible multimodal classification pipeline comparing classical
        machine learning baselines with specialized transformer architectures
        (BanglaBERT and XLM-RoBERTa).
  \item Coordinated campaign detection modules integrating semantic
        representations, temporal posting patterns, and network-level interaction
        graphs.
\end{enumerate}

\clearpage
\tableofcontents
\clearpage
\listoffigures
\clearpage
\listoftables
\clearpage
\pagenumbering{arabic}

% =====================================================================
%                            CHAPTER 1
% =====================================================================
\chapter{Introduction}
\label{ch:intro}

This chapter presents an overview of the proposed research on coordinated
propaganda detection in Bangla digital media. It outlines the project background,
problem statement, core motivation, research objectives, methodology, expected
outcomes, and the structural organization of the report.

\section{Project Overview}
Bangladesh's digital media ecosystem has become a major source of information for
millions of people, alongside continued growth in internet access and social media
use. According to DataReportal's \emph{Digital 2025: Bangladesh}, the country had
\textbf{77.7 million internet users} and \textbf{60.0 million social media user
identities} in early 2025, demonstrating how strongly online platforms now shape
news circulation and public opinion.

Within this environment, \textbf{propaganda} refers to strategically crafted
information designed to influence public attitudes, beliefs, or actions. In the
digital age, such content is often spread through coordinated campaigns, repeated
messaging, and networked amplification, making it more scalable and harder to
detect than isolated misinformation. This issue is especially important in
Bangladesh, where online platforms are deeply embedded in political communication,
election-related discourse, and reactions to controversial events, all of which
increase the potential impact of coordinated influence operations~\cite{hristakieva2022}.

\subsection*{Problem Statement}
Most automated propaganda detection systems to date have been developed for English
and other languages with abundant digital resources. Bangla, by contrast, still
lacks large annotated datasets and detection frameworks capable of identifying
coordinated propaganda across text, images, and cross-platform activity, which
leaves a substantial research and practical gap~\cite{hossain2020}. This project
responds to that gap by proposing a Bangla propaganda dataset and an AI-based
detection framework tailored to Bangladesh's online media environment. The goal is
to support more accurate, scalable, and context-sensitive detection of coordinated
propaganda campaigns in Bangla-language digital spaces~\cite{hossain2020,kamruzzaman2023}.

\section{Motivation}
The motivation for this project extends beyond building a detection system. In an
increasingly digital society, people rely on online information to understand
current events, form opinions, and make decisions that affect their communities and
democratic processes. When coordinated propaganda manipulates information at scale,
it can undermine public trust, distort discourse, and reduce people's ability to
make informed judgments~\cite{hristakieva2022}.

This project is grounded in the belief that individuals should be able to think
independently rather than have their perceptions shaped by coordinated
manipulation. Supporting that independence requires a digital environment in which
misleading influence campaigns can be identified and examined more effectively.
Technology alone cannot eliminate propaganda, but it can provide tools that improve
transparency and help users evaluate online content more critically~\cite{ng2023}.

From a research standpoint, the scarcity of Bangla-language resources for this task
presents an opportunity. Existing Bangla work has already shown the value of
annotated datasets and benchmark systems for low-resource language processing, and
the same need extends to multimodal and propaganda-oriented settings. In this way,
the project aims to contribute both to advances in AI research and to a more
trustworthy online information ecosystem in Bangladesh~\cite{hossain2020}.

\section{Objectives}
The primary purpose of this research is to build a foundational AI framework and
dataset for the Bangla digital media environment that can effectively detect and
analyze coordinated propaganda campaigns~\cite{ng2023}.

As this research is currently in its initial phase, the specific objectives are
enumerated below:
\begin{enumerate}
  \item \textbf{Create a large-scale, multimodal Bangla dataset:} Compile and
        annotate a comprehensive dataset spanning diverse domains, such as politics
        and news, including both textual and visual elements, using hybrid
        annotation pipelines that combine human expertise with LLM assistance to
        scale the process efficiently~\cite{hossain2020}.
  \item \textbf{Develop an AI model for propaganda detection:} Architect and train
        a multimodal machine learning framework optimized for the linguistic
        nuances of Bangla to classify manipulative and propagandistic
        content~\cite{hossain2020}.
  \item \textbf{Analyze coordinated campaigns:} Combine semantic, temporal, and
        network-level signals to detect organized, scalable influence operations
        rather than isolated pieces of misinformation~\cite{ng2023}.
  \item \textbf{Evaluate model performance:} Conduct baseline evaluations against
        current state-of-the-art models, such as fine-tuned transformers and
        cross-lingual models, to assess accuracy, scalability, and
        context-sensitivity in the Bangladeshi digital landscape~\cite{hossain2020}.
  \item \textbf{Establish open resources for future research:} Make the annotated
        dataset and architectural findings accessible to the academic community to
        support further advances in NLP, multimodal machine learning, and digital
        media analysis for low-resource languages~\cite{hossain2020}.
\end{enumerate}

\section{Methodology}
\label{sec:method-overview}
To achieve these objectives, the research follows a structured, multi-phase
methodology. Figure~\ref{fig:workflow} presents the complete pipeline at a glance:
raw Bangla content is ingested and normalized, passed through a hybrid
LLM-plus-human annotation stage to yield a verified multimodal dataset, and then
processed by three parallel analysis streams (text, visual, and network) whose
outputs are combined by a cross-modal campaign-fusion head and delivered through a
web dashboard and monitoring API.
\begin{figure}[H]
\centering
\begin{adjustbox}{max width=\textwidth,max height=0.85\textheight}
\begin{tikzpicture}[
  wfn/.style={rectangle,draw=black!70,rounded corners=2pt,align=center,
              font=\small,minimum height=13mm,text width=34mm,inner sep=4pt},
  wfsrc/.style ={wfn,fill=gray!10},
  wfproc/.style={wfn,fill=blue!7,minimum height=20mm},
  wfsto/.style ={wfn,fill=green!8},
  wfout/.style ={wfn,fill=orange!10,text width=55mm},
  stl/.style={rotate=90,anchor=south,font=\scriptsize\bfseries,align=center,text width=30mm},
  wfgrp/.style={draw=black!45,dashed,rounded corners=4pt,inner sep=8pt}
]
% ---- Stage 1
\node[wfsrc] (src)   at (-4.6,0)    {\textbf{Bangla Digital Media Sources}\\[2pt]{\scriptsize Facebook \quad Twitter/X \quad News Outlets}};
\node[wfsrc] (clean) at (0,0)       {\textbf{Data Cleaning}\\[2pt]{\scriptsize Unicode Normalization}};
\node[wfsrc] (annot) at (4.6,0)     {\textbf{Hybrid Annotation Pipeline}\\[2pt]{\scriptsize LLM Pre-Annotation\\+ Human Verification}};
\node[wfsto] (ds)    at (0,-2.4)    {\textbf{Verified Multimodal Dataset}};
\draw[flow] (src) -- (clean);
\draw[flow] (clean) -- (annot);
\draw[flow] (annot.south) |- (ds.east);
\node[wfgrp,fit=(src)(clean)(annot)(ds)] (st1) {};
\node[stl] at (st1.west) {1.\ Ingestion \& Preprocessing};

% ---- Stage 2
\node[wfproc] (txt) at (-4.6,-6.2) {\textbf{Text Stream}\\[2pt]{\scriptsize BanglaBERT / XLM-R\\Semantic \& Persuasive Cue Modeling}};
\node[wfproc] (vis) at (0,-6.2)    {\textbf{Visual Stream}\\[2pt]{\scriptsize ResNet / ViT\\Image--Text Alignment \& Meme Analysis}};
\node[wfproc] (net) at (4.6,-6.2)  {\textbf{Network Stream}\\[2pt]{\scriptsize Temporal Clustering\\Coordinated Account Interaction Graphs}};
\node[wfgrp,fit=(txt)(vis)(net)] (st2) {};
\node[stl] at (st2.west) {2.\ Multimodal AI Framework};

% bus: dataset -> three streams
\draw[line] (ds.south) -- (0,-4.4);
\draw[line] (-4.6,-4.4) -- (4.6,-4.4);
\draw[flow] (-4.6,-4.4) -- (txt.north);
\draw[flow] (0,-4.4)    -- (vis.north);
\draw[flow] (4.6,-4.4)  -- (net.north);

% ---- Stage 3
\node[wfout] (fuse) at (0,-9.8)  {\textbf{Cross-Modal Campaign\\Fusion Head}};
\node[wfout] (out)  at (0,-11.7) {\textbf{Propaganda Category}\\{\scriptsize Confidence \& Technique}};
\node[wfout] (dash) at (0,-13.6) {\textbf{Web Dashboard}\\{\scriptsize \& Monitoring API}};
\draw[flow] (fuse) -- (out);
\draw[flow] (out) -- (dash);
\node[wfgrp,fit=(fuse)(out)(dash)] (st3) {};
\node[stl] at (st3.west) {3.\ Synthesis \& User Deployment};

% bus: three streams -> fusion
\draw[line] (txt.south) -- (-4.6,-8.4) -- (0,-8.4);
\draw[line] (net.south) -- (4.6,-8.4)  -- (0,-8.4);
\draw[flow] (vis.south) -- (fuse.north);
\end{tikzpicture}
\end{adjustbox}
\caption{Overall workflow of the proposed Bangla propaganda detection framework.}
\label{fig:workflow}
\end{figure}

The workflow proceeds through five core phases~\cite{ng2023}:
\begin{itemize}
  \item \textbf{Data Collection and Integration:} Raw data containing text and
        visual elements across politics, sports, and entertainment is
        systematically gathered using automated web scrapers and validated
        manually.
  \item \textbf{Data Cleaning and Preprocessing:} Unstructured text is cleaned,
        standardized, tokenized, and normalized into machine-readable formats
        tailored for Bangla NLP tasks.
  \item \textbf{Dataset Annotation and Construction:} A hybrid pipeline uses LLMs
        for initial pre-annotation and span extraction, followed by human expert
        verification to ensure label quality and reduce
        bias~\cite{hossain2020}.
  \item \textbf{Multimodal Model Development:} Semantic features, temporal metrics,
        and social network structure are fused to pinpoint coordinated campaigns
        rather than standalone posts~\cite{ng2023}.
  \item \textbf{Evaluation and Comparative Analysis:} Benchmarks are run against
        transformer architectures and cross-lingual models to measure real-world
        performance on Bangladeshi digital media contexts~\cite{hossain2020}.
\end{itemize}

\section{Project Outcome}
By the end of this phase of the research, the project is expected to deliver the
following major concrete outputs~\cite{hossain2020}:
\begin{itemize}
  \item An annotated multimodal Bangla propaganda dataset covering multiple public
        domains.
  \item A trained AI framework designed for detecting coordinated digital
        propaganda.
  \item A comprehensive baseline evaluation comparing the proposed architecture
        against state-of-the-art models.
  \item Openly available dataset resources and benchmarks to foster future NLP and
        multimodal research for under-resourced languages~\cite{hristakieva2022}.
\end{itemize}

\section{Organization of the Report}
The remainder of this report is organized as follows:
\begin{itemize}
  \item \textbf{Chapter~\ref{ch:background} (Background):} Presents the necessary
        preliminaries, reviews similar applications and related research, and
        identifies the research gaps relevant to coordinated propaganda detection
        in Bangla online media.
  \item \textbf{Chapter~\ref{ch:design} (Project Design):} Describes the project
        requirements, context diagram, data flow diagrams, user interface design,
        detailed methodology, project plan, and task allocation.
  \item \textbf{Chapter~\ref{ch:standards} (Standards and Design Constraints):}
        Discusses the standards followed, realistic design constraints, cost
        analysis, and the mapping of the project to complex engineering problem
        solving and activity criteria.
\end{itemize}

% =====================================================================
%                            CHAPTER 2
% =====================================================================
\chapter{Background}
\label{ch:background}

This chapter provides the background knowledge needed to understand the proposed
work and discusses existing applications and research related to propaganda,
misinformation, and harmful-content detection. It also identifies the main research
gaps that motivate the proposed system.

\section{Preliminaries}
Social media has become an important platform for sharing news, opinions, images,
and other information. At the same time, it can also be used to spread propaganda,
misinformation, hate speech, manipulation, and other harmful content. Because
harmful information can spread very quickly, automatic detection systems are
becoming increasingly important.

\subsection{Propaganda}
Propaganda is information designed to influence people's opinions or behavior, often
by using selective information, emotional language, misleading presentation, or
other persuasive techniques. In social media, propaganda can appear as text,
images, videos, hashtags, or coordinated posts.

\subsection{Misinformation and Harmful Content}
Misinformation refers to incorrect or misleading information. Harmful online
content is a broader term that can include hate speech, offensive language,
harassment, manipulated information, and other content that may cause social harm.
These types of content can be difficult to detect automatically because their
meaning often depends on context.

\subsection{Machine Learning for Content Detection}
Machine learning (ML) allows a computer system to learn patterns from previously
labeled examples. Traditional ML methods such as Support Vector Machines (SVM) have
been widely used for text classification. More recent systems use deep learning and
transformer-based language models to understand more complex language
patterns~\cite{duridi2025,plikynas2025}.

\subsection{Deep Learning and Transformer Models}
Deep learning (DL) models can automatically learn useful features from data.
Transformer-based models are particularly useful for language tasks because they
can capture relationships between words over a long text. Models such as BERT and
language models adapted to specific languages have therefore become common in
harmful-content detection~\cite{duridi2025,lora2023,hasan2026}.

\subsection{Graph-Based Detection}
Graph-based approaches represent information as connected nodes and relationships.
They are useful when the relationship between words, users, posts, or accounts is
important. For example, propaganda can depend on relationships between different
parts of a text rather than on individual words alone. Ahmad et al.\ used a
hierarchical graph-based model to capture such relationships in propaganda
detection~\cite{ahmad2025}.

\subsection{Multimodal Detection}
Social media posts often contain more than one type of information. A post may
contain text and an image together. Multimodal detection combines these different
sources of information to make a prediction. Research has shown that image
information can provide important clues for fake-news and hateful-meme
detection~\cite{lin2024,hee2022}.

\subsection{Behavioral and Network Signals}
The behavior of a social-media account can also provide useful information.
Examples include how often an account posts, how it interacts with other users, and
whether it participates in coordinated activity. These signals can be useful when
the actual content is difficult to trust or is deliberately changed to avoid
detection~\cite{declerck2024,sela2025,schneider2026}.

\subsection{Low-Resource Languages}
A low-resource language is a language for which relatively few labeled datasets,
language models, or other computational resources are available. Bengali is an
example where researchers have identified limitations in annotated datasets and
language resources~\cite{lora2023,hasan2026,romim2022}. This makes harmful-content
detection more challenging than in high-resource languages such as English.

\section{Literature Review}
The reviewed literature focuses on computational methods for detecting propaganda,
misinformation, manipulation, hate speech, sarcasm, and related harmful content.
The studies cover traditional machine learning, deep learning, graph models,
transfer learning, multimodal methods, behavioral analysis, and large language
models (LLMs)~\cite{duridi2025,plikynas2025,ahmad2025,schneider2026}.

The literature can be divided into two areas: \textbf{similar applications} and
\textbf{related academic research}.

\subsection{Similar Applications}
Several existing tools and systems address parts of the problem of detecting
misinformation, manipulated information, or suspicious online activity. These
systems show that automated analysis can be provided through web-based interfaces
and APIs, although they do not necessarily perform exactly the same task as the
proposed system.
\begin{itemize}
  \item \textbf{Bot-detection and Coordination Analysis Systems:} Some systems
        focus on identifying automated or coordinated social-media accounts.
        Research-based systems such as those discussed by De~Clerck et al.\ analyze
        account behavior and coordination patterns to identify bot-like
        activity~\cite{declerck2024}. These systems are related to the proposed work
        because they show how user behavior can be used in addition to the content
        of a post.
  \item \textbf{Fact-checking Systems:} Fact-checking platforms and services allow
        users or applications to check whether claims have already been reviewed.
        Google provides Fact Check Tools that support searching and working with
        fact-check information through an API~\cite{google_factcheck}. Google also
        supports the \texttt{ClaimReview} format for representing fact-check
        results in a structured way~\cite{google_claimreview}. Such systems are
        related because they support verification of questionable information,
        although fact-checking is different from automatically detecting
        propaganda techniques.
  \item \textbf{Social-media Misinformation Analysis:} Other applications and
        research systems analyze how information spreads through social networks.
        The work of Sela et al.\ shows that repetition patterns and network
        structures can provide signals for identifying political influence in
        information spread~\cite{sela2025}. These systems are relevant because they
        demonstrate that harmful information can be analyzed at both the content
        level and the network level.
  \item \textbf{Limitations of Existing Applications:} A common limitation of
        existing applications is that they often focus on one specific problem,
        such as fact checking, bot detection, or misinformation analysis. A system
        that combines these ideas with low-resource language and multimodal
        analysis can address a broader set of challenges.
\end{itemize}

\subsection{Related Research}

\subsubsection*{Propaganda Detection Using Machine Learning and Deep Learning}
Traditional machine learning remains useful for some propaganda detection tasks.
Duridi et al.\ found that SVM performed better than AraBERT for propaganda
detection in their Arabic social-media dataset, while AraBERT performed better for
bias detection~\cite{duridi2025}. This result shows that the most advanced model is
not always the best model for every task.

Ahmad et al.\ proposed a Hierarchical Graph-based Integration Network (H-GIN) for
propaganda detection. The model uses graph-based information to capture long-range
and non-adjacent relationships that can be missed by standard sequence models. The
study reported an accuracy of 82\% for its evaluated task~\cite{ahmad2025}.

These studies show that both traditional and advanced models can be useful. The
choice of model depends on the task, the dataset, and the type of information that
needs to be captured.

\subsubsection*{Behavioral and Network-Based Detection}
Recent research has moved beyond text-only detection. De~Clerck et al.\ studied
coordinated and bot-like behavior on Twitter and showed that patterns of account
activity can help identify suspicious behavior~\cite{declerck2024}.

Sela et al.\ also found that repetition patterns and network structure can provide
useful signals for identifying political influence in information
spread~\cite{sela2025}.

Schneider et al.\ further developed this idea by focusing on user behavior rather
than only the content of posts. In their Reddit study, behavioral-policy
representations achieved a macro-F1 of 94.9\%, compared with 91.2\% for text
embeddings in the evaluated actor-detection task~\cite{schneider2026}.

These studies suggest that user behavior and network structure can provide
information that cannot be obtained from text alone.

\subsubsection*{Low-Resource Language Detection}
Low-resource languages face important challenges because of limited data and
language resources. Haider et al.\ studied manipulation detection in low-resource
languages and showed that transfer learning can help when only a small amount of
training data is available in the target language~\cite{haider2020}.

For Bengali, Lora et al.\ developed a transformer-based model for sarcasm
detection. Their results differed across datasets, showing that model performance
can depend strongly on the characteristics of the available corpus~\cite{lora2023}.

Romim et al.\ introduced the BD-SHS benchmark dataset for Bangla hate-speech
detection. The dataset contains manually annotated comments from different social
contexts, and the study reported a best F1-score of 91.0\% for its evaluated
system~\cite{romim2022}.

More recent work by Hasan et al.\ introduced BanglaMultiHate, which considers
several aspects of hate speech, including type, severity, and target. The study
evaluated classical machine learning, BanglaBERT, and LLM-based
approaches~\cite{hasan2026}.

Together, these studies show that successful moderation in Bangla requires larger
datasets, better annotation, and models that understand local language and social
context.

\subsubsection*{Multimodal Detection}
Lin et al.\ proposed a text-image fusion model for fake-news detection. Their
results demonstrate that combining text and image information can improve detection
compared with using only one modality~\cite{lin2024}.

Hee et al.\ studied multimodal hateful memes and found that images can provide
important information for classification. However, they also observed that
multimodal systems may learn unwanted biases and produce false
positives~\cite{hee2022}.

Wang et al.\ investigated image-based political propaganda and showed that visual
elements such as colors can be used in different ways by political
groups~\cite{wang2022}.

These findings suggest that text-only systems may miss important visual
information. Therefore, multimodal analysis is an important direction for future
harmful-content detection.

\subsubsection*{Large Language Models}
Large language models are a newer research direction in propaganda detection.
Pi\~na-Garc\'ia used few-shot LLaMA~3.2 to classify propagandistic tweets and
examine linguistic and coordination patterns in Mexican election-related
data~\cite{pina2025}.

Gaeta et al.\ proposed an LLM-based system that detects propaganda and identifies
the specific parts of a text that contain propagandistic content~\cite{gaeta2025}.

However, LLM performance depends strongly on data quality. Maarouf et al.\ showed
that models trained with weak labels achieved an AUC of 64.03, whereas models
trained with high-quality human annotations achieved an AUC of 92.25. Prompt-based
learning with a smaller set of high-quality samples achieved an AUC of
80.27~\cite{maarouf2024}.

This indicates that better models alone are not enough. High-quality training and
evaluation data are also essential.

\subsection{Summary of Literature}
\label{sec:lit-summary}
Table~\ref{tab:litsummary} provides a systematic comparative summary of the
reviewed literature, outlining the core computational methods, benchmark datasets,
key findings, and critical limitations of prior work.

\begingroup
\small
\setlength{\tabcolsep}{4pt}
\renewcommand{\arraystretch}{1.25}
\begin{longtable}{@{}L{3.1cm}C{0.8cm}L{2.2cm}L{3.5cm}L{3.3cm}@{}}
\caption{Summary of related research papers.}\label{tab:litsummary}\\
\toprule
\textbf{Paper Title \& Citation} & \textbf{Year} & \textbf{Dataset Used} & \textbf{Key Methodology \& Findings} & \textbf{Critical Limitations}\\
\midrule
\endfirsthead
\multicolumn{5}{@{}l}{\small\itshape Table~\thetable\ (continued)}\\
\toprule
\textbf{Paper Title \& Citation} & \textbf{Year} & \textbf{Dataset Used} & \textbf{Key Methodology \& Findings} & \textbf{Critical Limitations}\\
\midrule
\endhead
\midrule
\multicolumn{5}{r@{}}{\small\itshape Continued on next page}\\
\endfoot
\bottomrule
\endlastfoot

\textbf{Spread of propaganda by coordinated communities on social media}~\cite{hristakieva2022} & 2022 & Twitter (Italian elections \& COVID-19) &
Analyzed coordinated link-sharing networks and bot behavior using network graphs. Found that propaganda spreaders exhibit significantly higher coordination than benign users. &
Evaluated only on high-resource languages (Italian/English); text semantics were secondary to network topology; lacks multimodal analysis.\\
\addlinespace
\textbf{BanFakeNews: A dataset for detecting fake news in Bangla}~\cite{hossain2020} & 2020 & BanFakeNews (50,000 Bangla news articles) &
Landmark benchmark Bangla dataset for fake news detection; evaluated SVM, Random Forest, and LSTM architectures. &
Focuses strictly on factual veracity (fake vs.\ real) rather than rhetorical propaganda techniques; purely text-based without social context or coordination tracking.\\
\addlinespace
\textbf{BanMANI: Manipulated social media news in Bangla}~\cite{kamruzzaman2023} & 2023 & BanMANI (social media news posts) &
Categorizes social media news manipulation into multiple classes; demonstrates the efficacy of transfer learning. &
Small dataset size; does not analyze coordinated amplification campaigns or image--text multimodal relationships.\\
\addlinespace
\textbf{Coordinated information campaigns on social media}~\cite{ng2023} & 2023 & Twitter / Reddit multi-domain data &
Multi-faceted framework fusing semantic similarity, temporal proximity, and network clustering to detect coordinated operations. &
Computationally intensive; designed primarily for English; relies on extensive historical user activity graphs often unavailable in low-resource contexts.\\
\addlinespace
\textbf{Propaganda and bias in social media: Israel--Gaza war}~\cite{duridi2025} & 2025 & Arabic X (Twitter) posts &
Evaluated classical ML (SVM) against deep learning (AraBERT); SVM achieved superior performance on propaganda classification while AraBERT excelled at bias detection. &
Restricted to short text posts; does not address multimodal memes or cross-platform campaign coordination.\\
\addlinespace
\textbf{Hierarchical graph-based integration network for propaganda detection}~\cite{ahmad2025} & 2025 & SemEval propaganda corpus (English) &
H-GIN model capturing long-range, non-adjacent token relationships via graph neural networks; reported 82\% accuracy. &
Operates solely on English long-form news articles; high graph construction latency; no support for low-resource or code-mixed text.\\
\addlinespace
\textbf{Text-image multimodal fusion for fake news detection}~\cite{lin2024} & 2024 & Weibo \& Twitter multimodal benchmarks &
Cross-modal attention network fusing visual features (ResNet/ViT) with textual embeddings (BERT); showed multimodal fusion outperforms unimodal baselines. &
Evaluated on Chinese and English benchmarks; vulnerable to false positives when meme text and images are intentionally incongruent.\\
\addlinespace
\textbf{Explaining multimodal hateful meme detection models}~\cite{hee2022} & 2022 & Facebook Hateful Memes dataset &
Investigates multimodal interpretability; reveals that multimodal models frequently latch onto background image biases and surface textual heuristics. &
Demonstrates explainability challenges; lacks application to political propaganda in South Asian contexts.\\
\addlinespace
\textbf{Beyond content: Behavioral policies reveal actors in information operations}~\cite{schneider2026} & 2026 & Twitter \& Reddit coordinated influence sets &
Models user behavioral policies (posting timing, activity patterns) reaching 94.9\% macro-F1, outperforming text embeddings (91.2\%). &
Requires rich temporal account metadata; unviable for anonymous or newly created ephemeral accounts.\\
\addlinespace
\textbf{BD-SHS: Benchmark dataset for online Bangla hate speech}~\cite{romim2022} & 2022 & BD-SHS (Bangla social comments) &
Diverse benchmark across multiple social categories; achieved 91.0\% F1-score using transformer fine-tuning. &
Focuses on abusive language rather than persuasive propaganda techniques; lacks coordination and image features.\\
\addlinespace
\textbf{LLM-based multi-task Bangla hate speech detection}~\cite{hasan2026} & 2026 & BanglaMultiHate dataset &
Multi-task framework evaluating classical ML, BanglaBERT, and few-shot LLMs across hate speech type, severity, and target. &
Shows that while LLMs are promising, domain-specialized models (BanglaBERT) often yield superior and more consistent F1 on local cultural nuances.\\
\addlinespace
\textbf{HQP: A human-annotated dataset for detecting online propaganda}~\cite{maarouf2024} & 2024 & HQP (English articles) &
Proves data quality impact: weak labels gave 64.03 AUC, whereas human-annotated data achieved 92.25 AUC. Prompt-learning with small verified sets achieved 80.27 AUC. &
English-only; underscores the vital necessity of human verification over pure LLM automated annotation.\\
\end{longtable}
\endgroup

\section{Gap Analysis}
The reviewed literature shows several important gaps that can be addressed by the
proposed work:

\subsection{Gap 1: Limited Support for Low-Resource Languages}
Many existing systems are developed and tested mainly on high-resource languages.
Research on Bengali and Bangla shows that limited annotated datasets and language
resources remain major problems~\cite{lora2023,hasan2026,romim2022}. More work is
needed on systems that can effectively handle low-resource languages.

\subsection{Gap 2: Heavy Dependence on Text}
Many detection systems focus mainly on textual information. However, social-media
content often includes images, and important information may appear in the
relationship between text and images~\cite{lin2024,hee2022}. This creates a need
for systems that can analyze multiple modalities together.

\subsection{Gap 3: Limited Use of Behavioral Information}
Traditional content-based systems may fail when users intentionally change or
disguise their language. Behavioral and network-based research shows that user
activity and coordination can provide additional
evidence~\cite{declerck2024,sela2025,schneider2026}. However, these signals are not
always integrated with text and multimodal models.

\subsection{Gap 4: Data Quality and Annotation Problems}
The quality of the training data has a strong effect on model performance. HQP
demonstrates that high-quality human annotation can produce much better results
than weak labeling~\cite{maarouf2024}. Bangla research also emphasizes the
importance of culturally grounded and socially diverse
datasets~\cite{hasan2026,romim2022}.

\subsection{Gap 5: Generalization Across Datasets and Domains}
A model can perform well on one dataset but perform differently on another dataset.
The reviewed literature therefore indicates the need for evaluation across multiple
datasets, domains, and languages rather than relying on a single
benchmark~\cite{ahmad2025}.

\subsection{Gap 6: Explainability}
Many machine-learning and deep-learning systems provide a final classification
without clearly explaining why a particular post was detected as harmful. This can
be especially problematic for propaganda, sarcasm, memes, and culturally sensitive
content. More interpretable systems are therefore
needed~\cite{ahmad2025,hee2022,maarouf2024}.

\subsection{Proposed Research Opportunity}
Based on these gaps, there is an opportunity to develop a system that combines
multiple types of evidence. Such a system could use text analysis together with
image information and, where available, behavioral or network signals. For a
low-resource language such as Bangla, the system could also use language-specific
pretrained models and carefully annotated local datasets.

The main focus should therefore be on building a detection approach that is not
dependent on a single signal or a single dataset. It should be evaluated for
accuracy, robustness, and generalization across different types of online content.

\section{Summary}
This chapter presented the background concepts needed to understand propaganda and
harmful-content detection. It discussed machine learning, deep learning,
graph-based approaches, multimodal analysis, behavioral signals, and the challenges
of low-resource languages.

The literature review showed that researchers are moving from text-only systems
toward approaches that combine content, user behavior, network structure, images,
and large language models~\cite{duridi2025,plikynas2025,lora2023,hasan2026,ahmad2025,lin2024,hee2022,declerck2024,sela2025,schneider2026,romim2022,haider2020,wang2022,pina2025,gaeta2025,maarouf2024}.
The studies also show that no single method works equally well for every task.
Model performance depends on the problem, the dataset, the available language
resources, and the quality of annotation.

The gap analysis identified several areas that still need improvement, especially
low-resource language support, multimodal analysis, behavioral information, data
quality, cross-dataset evaluation, and explainability. These gaps provide the
motivation for developing a more integrated detection system that can work with
low-resource and multimodal social-media content.

% =====================================================================
%                            CHAPTER 3
% =====================================================================
\chapter{Project Design}
\label{ch:design}

\section{Project Design Overview}
This chapter presents the design of the proposed Bangla propaganda detection
system. It describes the system requirements, overall interaction and data flow,
user-interface design, the detailed methodology (including alternative solutions
considered and the reasons for selecting the proposed approach), the project plan,
and the allocation of tasks among team members.

\section{Requirement Analysis}
Requirement analysis identifies what the proposed system should do and what
constraints it should satisfy. The requirements are divided into functional and
nonfunctional requirements, followed by diagrams showing how users, data, and
system components interact.

\subsection{Functional and Nonfunctional Requirements}

\subsubsection*{Functional Requirements}
\begin{itemize}
  \item \textbf{FR1: User Input.} The system shall allow a user to enter or submit
        Bangla text for analysis.
  \item \textbf{FR2: Text Preprocessing.} The system shall clean and prepare the
        submitted Bangla text before sending it to the detection model.
  \item \textbf{FR3: Propaganda Detection.} The system shall analyze the input text
        and determine whether propaganda is present.
  \item \textbf{FR4: Propaganda Technique Identification.} When propaganda is
        detected, the system should identify the possible propaganda technique,
        such as loaded language, fear appeal, or bandwagon, based on the final
        label set used in the dataset.
  \item \textbf{FR5: Result Display.} The system shall display the prediction
        clearly to the user (Input: Bangla social-media text; Prediction:
        Propaganda Detected; Technique: Loaded Language).
  \item \textbf{FR6: Confidence Information.} The system may display a confidence
        score or confidence level associated with the prediction.
  \item \textbf{FR7: Dataset Management.} The research component of the system
        shall support the collection, cleaning, annotation, verification, and
        storage of Bangla propaganda examples.
  \item \textbf{FR8: LLM-Assisted Pre-Annotation.} The dataset-development
        pipeline shall allow an LLM to generate an initial label suggestion for
        collected Bangla content.
  \item \textbf{FR9: Human Verification.} The system or research workflow shall
        allow a human annotator to review and correct the LLM's suggested label.
  \item \textbf{FR10: Model Training.} The research pipeline shall support training
        and fine-tuning of the selected propaganda detection model using the
        human-verified dataset.
  \item \textbf{FR11: Model Evaluation.} The system shall calculate evaluation
        measures including accuracy, precision, recall, and F1-score.
  \item \textbf{FR12: Iterative Dataset Improvement.} The research workflow shall
        identify poorly performing classes and allow additional data to be
        collected or generated and verified.
\end{itemize}

\subsubsection*{Nonfunctional Requirements}
\begin{itemize}
  \item \textbf{NFR1: Usability.} The user interface should be simple enough that a
        user can submit Bangla text and understand the prediction without
        technical knowledge.
  \item \textbf{NFR2: Performance.} The system should provide predictions within a
        reasonable response time for normal user requests.
  \item \textbf{NFR3: Reliability.} The system should provide consistent results
        when the same input is analyzed under the same model configuration.
  \item \textbf{NFR4: Accuracy.} The model should achieve acceptable performance
        according to the selected evaluation metrics.
  \item \textbf{NFR5: Maintainability.} The system should be designed so that the
        detection model and dataset can be updated without rebuilding the complete
        application.
  \item \textbf{NFR6: Scalability.} The architecture should allow the dataset and
        model to grow beyond the initial 500--1,000 custom samples.
  \item \textbf{NFR7: Security.} User-submitted data and research data should be
        handled securely and unnecessary personal information should not be stored.
  \item \textbf{NFR8: Explainability.} Where feasible, the system should provide
        useful information about why a text was classified as potentially
        propagandistic rather than displaying only a class label.
  \item \textbf{NFR9: Language Support.} The main detection workflow shall support
        Bangla text.
\end{itemize}

% ---------------------------------------------------------------------
\subsection{Context Diagram}
The context diagram represents the proposed system as a single high-level process
and shows its relationship with external entities. The main external entities are:
\begin{enumerate}
  \item \textbf{User:} \emph{User $\rightarrow$ System:} Bangla text.
        \emph{System $\rightarrow$ User:} detection result, propaganda category,
        and optional confidence information.
  \item \textbf{Researcher / Administrator:} manages the research dataset and
        model-development process. \emph{Researcher $\rightarrow$ System:} raw
        data, annotation decisions, training configuration.
        \emph{System $\rightarrow$ Researcher:} dataset information, model
        results, evaluation metrics.
  \item \textbf{LLM Service:} used during dataset development as a pre-annotation
        assistant. \emph{System $\rightarrow$ LLM:} Bangla sample and annotation
        prompt. \emph{LLM $\rightarrow$ System:} suggested label, possible
        technique, and explanation. (The LLM output is \emph{not} considered the
        final ground truth; human verification is strictly required.)
  \item \textbf{External Data Sources:} publicly available datasets and other
        permitted sources provide the initial and additional Bangla data used
        during development. \emph{External Data Sources $\rightarrow$ System:}
        Bangla text/images and associated metadata where available.
\end{enumerate}

\begin{figure}[H]
\centering
\begin{adjustbox}{max width=\textwidth}
\begin{tikzpicture}
% central system
\node[draw=boxdraw,thick,fill=boxfill,rounded corners=6pt,minimum width=32mm,
      minimum height=66mm,text width=26mm,align=center,font=\footnotesize\bfseries] (S) at (0,0)
      {Bangla Propaganda Detection System};
% external entities
\node[ext] (U) at (-5.6, 2.2) {User};
\node[ext] (R) at (-5.6,-2.2) {Researcher / Administrator};
\node[ext] (L) at ( 5.6, 2.2) {LLM Service};
\node[ext] (E) at ( 5.6,-2.2) {External Data Sources};

% User <-> System
\draw[flow] (-4.29, 2.65) -- node[lab,above,text width=25mm]{Bangla text} (-1.6, 2.65);
\draw[flow] (-1.6, 1.75) -- node[lab,below,text width=25mm]{Detection result, category \& confidence} (-4.29, 1.75);
% Researcher <-> System
\draw[flow] (-4.29,-1.75) -- node[lab,above,text width=25mm]{Raw data, annotation decisions, training config} (-1.6,-1.75);
\draw[flow] (-1.6,-2.65) -- node[lab,below,text width=25mm]{Dataset info, model results, metrics} (-4.29,-2.65);
% System <-> LLM
\draw[flow] ( 1.6, 2.65) -- node[lab,above,text width=25mm]{Cleaned sample \& annotation prompt} ( 4.29, 2.65);
\draw[flow] ( 4.29, 1.75) -- node[lab,below,text width=25mm]{Suggested label, technique \& explanation} ( 1.6, 1.75);
% External data -> System
\draw[flow] ( 4.29,-2.2) -- node[lab,above,text width=25mm]{Raw Bangla text/images \& metadata} ( 1.6,-2.2);
\end{tikzpicture}
\end{adjustbox}
\caption{Context diagram --- the proposed system acts as the central component
connecting the user, research workflow, external datasets, and the LLM-based
annotation assistant.}
\label{fig:context}
\end{figure}

% ---------------------------------------------------------------------
\subsection{Data Flow Diagram}
\label{sec:dfd}

\subsubsection{Level 1 Data Flow Diagram}
The Level~1 Data Flow Diagram (DFD) expands the main system into its major internal
processes and data stores.

\paragraph{Processes.}
\begin{itemize}
  \item \textbf{P1: Data Collection and Management.} Collects Bangla data from
        existing datasets and permitted online sources.
  \item \textbf{P2: Data Preprocessing.} Cleans, normalizes, and prepares the
        collected data.
  \item \textbf{P3: LLM Pre-Annotation.} Sends the cleaned samples to the LLM and
        receives possible propaganda labels and explanations.
  \item \textbf{P4: Human Verification.} Allows the researcher or annotator to
        review the LLM output and produce the final verified label.
  \item \textbf{P5: Model Training.} Uses the verified dataset to fine-tune the
        selected machine-learning model.
  \item \textbf{P6: Model Evaluation.} Evaluates the trained model using accuracy,
        precision, recall, F1-score, and class-wise performance.
  \item \textbf{P7: User Prediction.} Receives new Bangla text from the user,
        preprocesses it, sends it to the trained model, and returns the prediction.
\end{itemize}

\paragraph{Main data stores.}
\begin{itemize}
  \item \textbf{D1: Raw Dataset} --- collected Bangla data before processing.
  \item \textbf{D2: Processed Dataset} --- cleaned and prepared data.
  \item \textbf{D3: Human-Verified Dataset} --- final dataset with human-confirmed
        labels.
  \item \textbf{D4: Trained Model} --- the trained propaganda detection model.
  \item \textbf{D5: Evaluation Results} --- performance measures and experimental
        results.
\end{itemize}

\begin{figure}[H]
\centering
\begin{adjustbox}{max width=\textwidth,max height=0.84\textheight}
\begin{tikzpicture}
% external entities (left)
\node[ext]  (eds) at (-5.3,  0.0) {External Data Sources};
\node[ext]  (llm) at (-5.3, -4.8) {LLM Service};
\node[ext]  (res) at (-5.3, -7.2) {Researcher / Expert};
\node[ext]  (usr) at (-5.3,-15.9) {End User};
% processes (centre)
\node[proc,text width=30mm] (p1) at (0,  0.0) {P1: Data Collection \& Management};
\node[proc,text width=30mm] (p2) at (0, -2.4) {P2: Data Preprocessing};
\node[proc,text width=30mm] (p3) at (0, -4.8) {P3: LLM Pre-Annotation};
\node[proc,text width=30mm] (p4) at (0, -7.2) {P4: Human Verification};
\node[proc,text width=30mm] (p5) at (0, -9.6) {P5: Model Training};
\node[proc,text width=30mm] (p6) at (0,-12.9) {P6: Model Evaluation (uses test split of D3)};
\node[proc,text width=30mm] (p7) at (0,-15.9) {P7: User Prediction};
% data stores (right)
\node[store,text width=26mm] (d1) at (5.3, -1.2) {D1: Raw Dataset};
\node[store,text width=26mm] (d2) at (5.3, -3.6) {D2: Processed Dataset};
\node[store,text width=26mm] (d3) at (5.3, -8.4) {D3: Human-Verified Dataset};
\node[store,text width=26mm] (d4) at (5.3,-11.0) {D4: Trained Model (weights \& config)};
\node[store,text width=26mm] (d5) at (5.3,-12.9) {D5: Evaluation Results};
\node[ext]   (r2) at (5.3,-14.9) {Researcher / Expert};

% ---- flows
\draw[flow] (eds) -- node[lab,above]{Raw content} (p1);
\draw[flow] (p1.east) -| (d1.north);
\draw[flow] (d1.west) -| (p2.north);
\draw[flow] (p2.east) -| (d2.north);
\draw[flow] (d2.west) -| (p3.north);
\draw[flow,<->] (p3.west) -- node[lab,above,text width=20mm]{Prompts \& suggestions} (llm.east);
\draw[flow] (p3.south) -- node[lab,right]{Pre-annotated} (p4.north);
\draw[flow,<->] (p4.west) -- node[lab,above,text width=20mm]{Review \& correct} (res.east);
\draw[flow] (p4.east) -| (d3.north);
\draw[flow] (d3.west) -| (p5.north);
\draw[flow] (p5.east) -| (d4.north);
\draw[flow] (d4.west) -| (p6.north);
\draw[flow] (p6.east) -- (d5.west);
\draw[flow] (d5.south) -- node[lab,right]{Results} (r2.north);
\draw[flow] (d4.east) -- ++(0.9,0) |- (p7.east);
\node[lab,above,text width=16mm] at (2.8,-15.9) {Inference weights};
% user <-> P7
\draw[flow] ($(usr.east)+(0,0.25)$) -- node[lab,above]{Bangla text} ($(p7.west)+(0,0.25)$);
\draw[flow] ($(p7.west)-(0,0.25)$) -- node[lab,below,text width=22mm]{Prediction, category \& explanation} ($(usr.east)-(0,0.25)$);
\end{tikzpicture}
\end{adjustbox}
\caption{Level~1 Data Flow Diagram (DFD) --- the dual-workflow architecture: the
research and model-development pipeline and the end-user prediction workflow.}
\label{fig:dfd1}
\end{figure}

\subsubsection{Detailed Level 2 Data Flow Diagram}
\label{sec:dfd2}
Figure~\ref{fig:dfd2} expands the Level~1 processes to show their internal
processing. Processes~1 and~2 harvest and normalize both text and images.
Processes~3 and~4 form the hybrid annotation engine, where few-shot prompts drive an
LLM to produce candidate suggestions that are placed in a buffer and then verified
by human annotators in a dedicated workspace before being written to the verified
ground-truth store (D3). Process~5 tokenizes verified text, encodes it with a
BanglaBERT/XLM-R backbone, encodes images with a ViT/ResNet extractor, and joins the
two through a cross-modal attention fusion layer feeding the propaganda technique
classification head. Processes~6 and~7 serve live user input through the inference
engine and pair each prediction with a feature-attribution explainability module.

\begin{figure}[H]
\centering
\begin{adjustbox}{max width=\textwidth,max height=0.86\textheight}
\begin{tikzpicture}
% ===== Processes 1 & 2
\node[blk]  (raw)   at (-5.25, 0.0) {Raw Online Bangla Content};
\node[proc] (scr)   at (-1.75, 0.0) {Automated Scraper \& API Harvester};
\node[proc] (img)   at ( 1.75, 0.0) {Image Preprocessing: Normalization \& Resizing};
\node[proc] (clean) at ( 1.75,-1.9) {Bangla Text Cleaner: Emoji, Punctuation, Normalizer};
\draw[grp] (-6.8,0.95) rectangle (6.8,-2.85);
\node[grptitle,anchor=south west] at (-6.8,0.95) {Processes 1 \& 2: Ingestion \& Preprocessing};
\draw[flow] (raw) -- (scr);
\draw[flow] (scr) -- (img);
\draw[flow] (scr.south) |- (clean.west);

% ===== Processes 3 & 4
\node[proc]  (prompt) at ( 1.75,-4.5) {Few-Shot Prompt Constructor};
\node[proc]  (llm)    at (-1.75,-4.5) {LLM Pre-Annotation API};
\node[store] (buf)    at (-5.25,-4.5) {Candidate Suggestion Buffer};
\node[proc]  (ui)     at (-5.25,-6.4) {Human Annotator Verification Workspace};
\node[store] (gt)     at (-1.75,-6.4) {D3: Verified Ground-Truth Store};
\draw[grp] (-6.8,-3.55) rectangle (6.8,-7.35);
\node[grptitle,anchor=south west] at (-6.8,-3.55) {Processes 3 \& 4: Hybrid Annotation Engine};
\draw[flow] (clean.south) -- (prompt.north);
\draw[flow] (prompt) -- (llm);
\draw[flow] (llm) -- (buf);
\draw[flow] (buf) -- (ui);
\draw[flow] (ui) -- (gt);

% ===== Process 5
\node[proc] (tok)  at (-1.75, -9.0) {Bangla WordPiece / SentencePiece Tokenizer};
\node[proc] (bb)   at ( 1.75, -9.0) {BanglaBERT / XLM-R Transformer Backbone};
\node[proc] (fuse) at ( 1.75,-10.9) {Cross-Modal Attention Fusion Layer};
\node[proc] (vit)  at ( 5.25,-10.9) {Vision Transformer / ResNet Visual Extractor};
\node[proc] (cls)  at (-1.75,-10.9) {Propaganda Technique Classification Head};
\draw[grp] (-6.8,-8.05) rectangle (6.8,-11.85);
\node[grptitle,anchor=south east] at (6.8,-8.05) {Process 5: Feature Extraction \& Model Training};
\draw[flow] (gt.south) -- (tok.north);
\draw[flow] (tok) -- (bb);
\draw[flow] (bb) -- (fuse);
\draw[flow] (vit) -- (fuse);
\draw[flow] (fuse) -- (cls);
% image path (outside right)
\draw[dflow] (img.east) -- (7.05,0) |- (vit.east);
\node[lab,rotate=90,anchor=north] at (7.1,-5.5) {Image features};

% ===== Processes 6 & 7
\node[ext]  (uin)  at (-5.25,-14.4) {User Bangla Text / URL};
\node[proc] (inf)  at (-1.75,-14.4) {Model Inference Engine};
\node[proc] (pout) at ( 1.75,-13.5) {Propaganda Label \& Technique Identification};
\node[proc] (shap) at ( 1.75,-15.3) {Feature Attribution \& Explainability Module};
\node[proc] (disp) at ( 5.25,-14.4) {Web Results Dashboard};
\draw[grp] (-6.8,-12.55) rectangle (6.8,-16.25);
\node[grptitle,anchor=south east] at (6.8,-12.55) {Processes 6 \& 7: Prediction \& Explainability};
\draw[flow] (uin) -- (inf);
\draw[flow] (inf.east)  -- ++(0.43,0) |- (pout.west);
\draw[flow] (inf.east)  -- ++(0.43,0) |- (shap.west);
\draw[flow] (pout.east) -- ++(0.43,0) |- (disp.west);
\draw[flow] (shap.east) -- ++(0.43,0) |- (disp.west);
\draw[dflow] (cls.south) -- (inf.north);
\node[lab,anchor=west] at (-1.6,-12.2) {Trained weights};
\end{tikzpicture}
\end{adjustbox}
\caption{Detailed Level~2 Data Flow Diagram (DFD) for the proposed system.}
\label{fig:dfd2}
\end{figure}

% ---------------------------------------------------------------------
\subsection{UI Design}
The proposed application uses a simple, intuitive interface so that users can
easily submit Bangla content and understand the result.

\subsubsection*{Main Screen}
The main screen contains the application title (\textbf{Bangla Propaganda Detection
System}), a short explanation of the system purpose, a Bangla text input area, an
``Analyze'' button, and a result section.

\begin{figure}[H]
\centering
\begin{adjustbox}{max width=\textwidth}
\begin{tikzpicture}[ui/.style={font=\footnotesize,anchor=west}]
\draw[boxdraw,thick,rounded corners=4pt,fill=white] (-6.2,0) rectangle (6.2,-9.4);
\node[font=\large\bfseries] at (0,-0.75) {Bangla Propaganda Detection System};
\node[font=\footnotesize] at (0,-1.4) {Analyze Bangla digital content for persuasive \& propaganda cues};
\draw[boxdraw] (-5.8,-1.85) -- (5.8,-1.85);

\node[ui,font=\footnotesize\bfseries] at (-5.8,-2.4) {Enter Bangla Text:};
\draw[gray,rounded corners=2pt,fill=gray!6] (-5.8,-2.75) rectangle (5.8,-4.55);
\node[ui,text=gray] at (-5.6,-3.15) {\bn{এখানে বাংলা সোশ্যাল মিডিয়া পোস্ট বা খবর লিখুন\ldots}};

\node[draw=boxdraw,fill=uiublue,text=white,rounded corners=3pt,minimum width=32mm,
      minimum height=8mm,font=\footnotesize\bfseries] at (0,-5.3) {Analyze};
\draw[boxdraw] (-5.8,-6.0) -- (5.8,-6.0);

\node[ui,font=\footnotesize\bfseries] at (-5.8,-6.5) {Result:};
\node[ui,anchor=north west,text width=11cm,align=left] at (-5.6,-6.9)
  {\textbf{Prediction:} Propaganda Detected\\[2pt]
   \textbf{Technique:} Loaded Language\\[2pt]
   \textbf{Confidence:} Medium\\[2pt]
   \textbf{Explanation:} The text contains emotionally strong wording that may influence the reader.};
\end{tikzpicture}
\end{adjustbox}
\caption{Main screen result interface.}
\label{fig:ui-main}
\end{figure}

\subsubsection*{Result Interface}
The result page clearly separates the input, prediction, and explanation:
\textbf{Prediction:} Propaganda Detected; \textbf{Possible Technique:} Loaded
Language; \textbf{Confidence:} Medium; \textbf{Explanation:} the text contains
emotionally strong wording that may influence the reader. The explanation is
treated as supporting information rather than as a replacement for the
classification result.

\subsubsection*{Annotation Interface}
A separate interface is used by the researcher or annotator to support the LLM
pre-annotation $\rightarrow$ human verification workflow.

\begin{figure}[H]
\centering
\begin{adjustbox}{max width=\textwidth}
\begin{tikzpicture}[ui/.style={font=\footnotesize,anchor=west},
                    b/.style={ui,font=\footnotesize\bfseries}]
\draw[boxdraw,thick,rounded corners=4pt,fill=white] (-6.2,0) rectangle (6.2,-11.0);
\node[font=\large\bfseries] at (0,-0.7) {Researcher / Annotator Interface};
\draw[boxdraw] (-5.8,-1.2) -- (5.8,-1.2);

\node[ui] at (-5.8,-1.7) {\textbf{Sample ID:} \#1042};
\node[ui] at (-5.8,-2.3) {\textbf{Text:} ``\ldots'' (Bangla sample text)};

\node[b] at (-5.8,-3.1) {LLM Suggested Label:};
\node[ui,draw=boxdraw,fill=boxfill,inner sep=3pt] at (-2.2,-3.1) {Propaganda Detected};
\node[b] at (-5.8,-3.8) {Suggested Technique:};
\node[ui,draw=boxdraw,fill=boxfill,inner sep=3pt] at (-2.2,-3.8) {Fear Appeal};
\node[b] at (-5.8,-4.5) {LLM Rationale:};
\node[ui] at (-2.2,-4.5) {Mentions severe imminent societal threat};

\node[b] at (-5.8,-5.4) {Human Decision:};
\draw (-5.5,-6.0) circle (0.11);  \node[ui] at (-5.25,-6.0) {Accept LLM Suggestion};
\draw (-5.5,-6.65) circle (0.11); \node[ui] (mod) at (-5.25,-6.65) {Modify Label:};
\node[ui,draw=gray,inner sep=3pt,right=3mm of mod] {Select Technique $\blacktriangledown$};
\draw (-5.5,-7.3) circle (0.11);  \node[ui] at (-5.25,-7.3) {Mark Ambiguous / Sarcastic};

\node[b] at (-5.8,-8.05) {Annotator Notes:};
\draw[gray,rounded corners=2pt,fill=gray!6] (-5.8,-8.35) rectangle (5.8,-9.35);

\node[draw=boxdraw,fill=uiublue,text=white,rounded corners=3pt,minimum width=48mm,
      minimum height=8mm,font=\footnotesize\bfseries] at (0,-10.2) {Confirm Ground Truth};
\end{tikzpicture}
\end{adjustbox}
\caption{Annotation interface.}
\label{fig:ui-annot}
\end{figure}

% ---------------------------------------------------------------------
\section{Detailed Methodology and Design}
\label{sec:detailed-method}
This section explains how the proposed system will be developed and why the
selected solutions are appropriate. Several alternative approaches were considered
before selecting the proposed combination of existing datasets, human-verified
custom data, transformer-based models, and LLM-assisted annotation.

\subsection{Overall Design Approach}
The project follows an iterative design rather than attempting to build the
complete system in a single step:
\[
\begin{gathered}
\text{Existing Dataset} \rightarrow \text{Baseline Pipeline} \rightarrow \text{Custom Dataset} \\
\rightarrow \text{Annotation} \rightarrow \text{Model Training} \rightarrow \text{Evaluation} \rightarrow \text{Improvement}
\end{gathered}
\]
This approach reduces development risk because the model-development pipeline can
be tested before the custom dataset is complete.

\subsection{Alternative Solutions Considered}

\subsubsection{Alternative 1: Building a Large Dataset Completely From Scratch}
One possible approach was to immediately collect and manually annotate a very large
Bangla dataset (e.g., attempting to label 10,000--50,000 samples before beginning
model development).
\begin{itemize}
  \item \textbf{Advantages:} Labels would be created specifically for the research
        task; the dataset could be designed around the exact propaganda categories
        required.
  \item \textbf{Disadvantages:} Manual annotation would require a large amount of
        time; annotation cost would increase significantly; difficult and
        ambiguous examples would require additional review; model development
        would need to wait until a large dataset was completed.
  \item \textbf{Selected Approach:} The proposed research instead begins with
        existing datasets and gradually creates a smaller custom dataset. This
        allows the technical pipeline to be developed immediately and reduces the
        initial annotation burden.
\end{itemize}

\subsubsection{Alternative 2: Fully Automatic LLM Annotation}
Another possible solution was to ask an LLM to label all collected samples and
directly use those labels as the dataset ground truth.
\begin{itemize}
  \item \textbf{Advantages:} Very fast compared with manual annotation; large
        quantities of data can be labeled automatically; reduces human workload.
  \item \textbf{Disadvantages:} LLM predictions can be incorrect; ambiguous Bangla
        language may be interpreted incorrectly; the LLM may introduce systematic
        labeling bias; using LLM labels as ground truth may reduce dataset
        reliability.
  \item \textbf{Selected Approach:} The proposed system uses the LLM only for
        pre-annotation. Every final label must be checked by a human annotator.
        Therefore
        \[\text{LLM Suggestion}\neq\text{Ground Truth},\]
        and the ground truth is the final human-verified label.
\end{itemize}

\subsubsection{Alternative 3: Complete Manual Annotation}
A third approach is to manually annotate every collected sample without using an LLM.
\begin{itemize}
  \item \textbf{Advantages:} Direct human judgment; better control over the
        annotation process; no dependence on LLM-generated suggestions.
  \item \textbf{Disadvantages:} Very time-consuming; more expensive in terms of
        researcher effort; difficult to scale as the dataset grows; annotators may
        become tired, which can reduce consistency.
  \item \textbf{Selected Approach:} The proposed LLM-assisted approach provides a
        compromise. The LLM performs an initial classification, while humans make
        the final decision. This reduces repetitive work without removing human
        control.
\end{itemize}

\subsubsection{Alternative 4: Traditional Machine Learning}
Traditional models such as SVM, Logistic Regression, or Naive Bayes could be used as
the main final classifier.
\begin{itemize}
  \item \textbf{Advantages:} Simple to implement; fast to train; computationally
        inexpensive; useful as baseline models.
  \item \textbf{Disadvantages:} Usually require manually designed text features;
        may have difficulty capturing complex contextual relationships in Bangla;
        may perform poorly when propaganda depends on wider context.
  \item \textbf{Selected Approach:} A traditional model is still implemented as a
        baseline. However, the main experiments investigate transformer-based
        models because they capture contextual language information more
        effectively.
\end{itemize}

\subsubsection{Alternative 5: BanglaBERT vs.\ XLM-R}
Two candidate transformer approaches are particularly relevant:
\begin{itemize}
  \item \textbf{BanglaBERT.} \emph{Advantages:} designed specifically for Bangla;
        suitable for Bangla text classification; can provide strong
        language-specific representations. \emph{Disadvantages:} primarily focused
        on Bangla; may be less suitable if the project is later expanded to
        multiple languages.
  \item \textbf{XLM-R.} \emph{Advantages:} supports many languages; can support
        multilingual expansion; useful if the system later needs to process both
        Bangla and other languages. \emph{Disadvantages:} not designed only for
        Bangla; a Bangla-specific model may be more suitable for some Bangla-only
        tasks.
  \item \textbf{Selected Approach:} BanglaBERT is used as the primary candidate and
        XLM-R as a comparison model. This allows the research to determine whether
        a Bangla-specific model or a multilingual model performs better on the
        custom dataset instead of assuming the answer beforehand.
\end{itemize}

\subsubsection{Alternative 6: Text-Only vs.\ Multimodal Detection}
A text-only system is simpler because it requires only Bangla text. However,
propaganda can also appear through images, memes, and combinations of text and
visual content. Existing research shows that multimodal information can contain
useful detection signals.
\begin{itemize}
  \item \textbf{Text-Only Approach.} \emph{Advantages:} easier to implement;
        requires fewer computational resources; easier dataset collection;
        suitable for the initial prototype. \emph{Disadvantages:} cannot analyze
        visual propaganda; may miss important information in images.
  \item \textbf{Multimodal Approach.} \emph{Advantages:} uses both text and image
        information; more suitable for posts containing memes or image-based
        messages. \emph{Disadvantages:} more complex; requires image data and
        suitable multimodal models; requires more computational resources;
        annotation becomes more difficult.
  \item \textbf{Selected Approach:} The initial system focuses on Bangla text
        detection so that the core pipeline can be developed and evaluated
        properly. Multimodal support is investigated as an extension where
        suitable image--text data are available. This follows the project's
        iterative strategy: first make the core system reliable, then increase its
        complexity.
\end{itemize}

\subsection{Design Rationale and Trade-off Evaluation}
Table~\ref{tab:tradeoff} consolidates the alternatives above into a single
evaluation matrix, stating for each design layer what was considered and why the
selected option was chosen.

\begingroup
\small
\renewcommand{\arraystretch}{1.3}
\begin{longtable}{@{}L{2.9cm}L{3.9cm}L{6.1cm}@{}}
\caption{Design rationale and trade-off evaluation matrix.}\label{tab:tradeoff}\\
\toprule
\textbf{Design Layer} & \textbf{Considered Alternative} & \textbf{Selected Choice \& Engineering Rationale}\\
\midrule
\endfirsthead
\multicolumn{3}{@{}l}{\small\itshape Table~\thetable\ (continued)}\\
\toprule
\textbf{Design Layer} & \textbf{Considered Alternative} & \textbf{Selected Choice \& Engineering Rationale}\\
\midrule
\endhead
\bottomrule
\endlastfoot

\textbf{Dataset Creation Strategy} & Fully manual annotation from scratch (10,000+ samples). &
\textbf{Hybrid LLM pre-annotation + human verification:} reduces manual annotator fatigue and labeling cost while preserving ground-truth fidelity through expert human sign-off~\cite{maarouf2024}.\\
\addlinespace
\textbf{Annotation Authority} & Fully autonomous LLM zero-shot labeling. &
\textbf{Human-in-the-loop ground truth:} pure LLM labels can suffer from hallucination and cultural blindness in Bangla idioms; human verification protects label quality~\cite{hasan2026,maarouf2024}.\\
\addlinespace
\textbf{Core NLP Backbone} & Classical machine learning (TF-IDF + SVM / Naive Bayes). &
\textbf{Fine-tuned pretrained transformers (BanglaBERT \& XLM-R):} capture bidirectional contextual semantics, sarcasm, and rhetorical nuance that $n$-gram models miss. Classical ML is retained strictly as the E1 baseline.\\
\addlinespace
\textbf{Monolingual vs.\ Multilingual} & Generic multilingual models only (mBERT / XLM-R). &
\textbf{Dual evaluation (BanglaBERT vs.\ XLM-R):} BanglaBERT provides vocabulary specialization for Bengali script, while XLM-R tests cross-lingual robustness against code-mixed Bangla--English text.\\
\addlinespace
\textbf{Modality Scope} & Full multimodal (video + audio + text + network) from day one. &
\textbf{Phased, extensible modular architecture:} Phase~1 establishes robust text detection and dataset verification; Phase~2 adds image--text meme fusion; Phase~3 integrates graph-based coordination metrics.\\
\addlinespace
\textbf{Class Imbalance Strategy} & Standard cross-entropy on the raw class distribution. &
\textbf{Cost-sensitive loss weighting + strategic oversampling:} prevents bias toward the dominant ``No Propaganda'' class and supports viable recall on minority techniques (e.g., Bandwagon, Fear Appeal).\\
\end{longtable}
\endgroup

\subsection{Dataset Design}
The dataset design follows a three-stage methodology:
\begin{itemize}
  \item \textbf{Stage 1: Existing Dataset.} Existing publicly available Bangla
        datasets are used to develop and test the initial pipeline.
  \item \textbf{Stage 2: Custom Dataset.} Approximately 500--1,000 Bangla samples
        will initially be collected from appropriate public sources.
  \item \textbf{Stage 3: Dataset Improvement.} After model evaluation,
        underrepresented classes will be identified and additional real examples
        will be collected where necessary. Synthetic augmentation may be
        investigated only when real examples remain insufficient, following the
        workflow
        \[\text{Generate}\rightarrow\text{Filter}\rightarrow\text{Human Check}\rightarrow\text{Training Use}.\]
        Synthetic examples will be stored separately from real-world examples
        during evaluation.
\end{itemize}

\subsection{Annotation Design}
The annotation system is based on a two-level process:
\begin{itemize}
  \item \textbf{Level 1: LLM Pre-Annotation.} The LLM receives a Bangla sample and
        suggests whether propaganda may be present, the possible propaganda
        technique, confidence or certainty, and a short explanation.
  \item \textbf{Level 2: Human Verification.} The human annotator checks the
        prediction and assigns the final label.
\end{itemize}
The final dataset therefore contains
\[\text{Original Text}+\text{LLM Suggestion}+\text{Human Decision}+\text{Final Label}.\]
This information can also help researchers analyze where the LLM performs well or
poorly.

\subsection{Class Imbalance Design}
The dataset may contain significantly different numbers of examples for different
propaganda categories. For example: \emph{No Propaganda}: 700; \emph{Loaded
Language}: 100; \emph{Fear Appeal}: 40; \emph{Bandwagon}: 20; \emph{Other}: 10.
The project will monitor class distributions after each data-collection cycle.
Possible solutions include:
\begin{enumerate}
  \item \textbf{Additional real-data collection:} more minority-class examples will
        be collected where possible.
  \item \textbf{Class weighting:} the training algorithm may assign greater
        importance to minority classes.
  \item \textbf{Oversampling:} minority examples may be sampled more frequently
        during training.
  \item \textbf{Synthetic augmentation:} LLM-generated examples may be investigated
        after human filtering.
\end{enumerate}
The chosen technique will be determined experimentally rather than assumed to be
effective in advance.

\subsection{Training and Testing Design}
\begin{itemize}
  \item The human-verified dataset will be divided into \textbf{training},
        \textbf{validation}, and \textbf{test} subsets.
  \item The training set will be used for model learning.
  \item The validation set will be used for model selection and parameter tuning.
  \item The test set will be kept separate and used for final evaluation.
  \item The test data must not be used during training or annotation decisions in a
        way that causes data leakage.
\end{itemize}

\subsection{Experimental Design}
The project will compare multiple approaches rather than reporting only one model.
The experimental structure is shown in Table~\ref{tab:experiments}.

\begin{table}[H]
\centering
\caption{Experimental design.}\label{tab:experiments}
\renewcommand{\arraystretch}{1.25}
\begin{tabular}{@{}cl@{}}
\toprule
\textbf{Experiment} & \textbf{Approach}\\
\midrule
E1 & Traditional ML baseline\\
E2 & BanglaBERT\\
E3 & XLM-R\\
E4 & BanglaBERT + class-imbalance strategy\\
E5 & XLM-R + class-imbalance strategy\\
E6 & Optional LLM-assisted / synthetic augmentation experiment\\
\bottomrule
\end{tabular}
\end{table}
The final selection will be based on measured performance rather than assuming that
the most complex model will perform best.

\subsection{Evaluation Design}
The main performance measures will be:
\begin{itemize}
  \item \textbf{Accuracy:} overall proportion of correctly classified examples.
  \item \textbf{Precision:} how many predicted positive examples are actually positive.
  \item \textbf{Recall:} how many actual positive examples are correctly detected.
  \item \textbf{F1-score:} the harmonic mean of precision and recall.
\end{itemize}
For an imbalanced multi-class dataset, \textbf{macro-F1} and \textbf{class-wise
F1-score} will receive particular attention because they provide a clearer view of
minority-class performance. A confusion matrix will also be used to identify which
propaganda categories are commonly confused with one another.

\subsection{Error Analysis}
After evaluation, incorrectly classified examples will be manually inspected. The
researchers will investigate whether errors are caused by ambiguous language,
sarcasm, insufficient training examples, class imbalance, cultural context,
code-mixing, similar propaganda categories, or limitations of the selected model.
The error analysis will guide the next iteration of dataset collection and model
training.

\subsection{Final System Architecture}
The final proposed architecture is divided into two connected parts: a research
pipeline responsible for building and improving the model, and a user application
that uses the trained model to provide predictions (Figure~\ref{fig:arch}).

\begin{figure}[H]
\centering
\begin{adjustbox}{max width=\textwidth}
\begin{tikzpicture}
% ===== Research pipeline
\node[blk]   (ds)  at (-5.1, 0.0) {Data Sources};
\node[blk]   (dc)  at (-1.7, 0.0) {Data Collection};
\node[blk]   (dp)  at ( 1.7, 0.0) {Pre-processing};
\node[blk]   (lpa) at ( 5.1, 0.0) {LLM Pre-Annotation};
\node[blk]   (hv)  at ( 5.1,-1.9) {Human Verification};
\node[store] (vd)  at ( 1.7,-1.9) {Verified Dataset};
\node[blk]   (mt)  at (-1.7,-1.9) {Model Training};
\node[blk]   (ev)  at (-1.7,-3.8) {Evaluation};
\node[blk]   (mi)  at (-5.1,-3.8) {Model Improvement};
\draw[grp] (-6.8,0.95) rectangle (6.8,-4.75);
\node[grptitle,anchor=south west] at (-6.8,0.95) {Research Pipeline};
\draw[flow] (ds) -- (dc);  \draw[flow] (dc) -- (dp);  \draw[flow] (dp) -- (lpa);
\draw[flow] (lpa) -- (hv); \draw[flow] (hv) -- (vd);  \draw[flow] (vd) -- (mt);
\draw[flow] (mt) -- (ev);  \draw[flow] (ev) -- (mi);
\draw[dflow] (mi.north) -- (-5.1,-0.95) -| (dc.south);
\node[lab,anchor=west] at (-5.0,-2.3) {iterate};

% ===== User application
\node[ext]  (u)   at ( 5.1,-7.5) {User};
\node[blk]  (bti) at ( 1.7,-7.5) {Bangla Text Input};
\node[blk]  (up)  at (-1.7,-7.5) {Pre-processing};
\node[proc] (tpm) at (-5.1,-7.5) {Trained Propaganda Model};
\node[blk]  (pr)  at (-5.1,-9.4) {Prediction};
\node[blk]  (tce) at (-1.7,-9.4) {Technique / Confidence / Explanation};
\node[blk]  (res) at ( 1.7,-9.4) {User Result Display};
\draw[grp] (-6.8,-6.55) rectangle (6.8,-10.35);
\node[grptitle,anchor=south east] at (6.8,-6.55) {User Application};
\draw[flow] (u) -- (bti);  \draw[flow] (bti) -- (up);  \draw[flow] (up) -- (tpm);
\draw[flow] (tpm) -- (pr); \draw[flow] (pr) -- (tce);  \draw[flow] (tce) -- (res);

% deploy link
\draw[dflow] (ev.south) -- (-1.7,-5.6) -| (tpm.north);
\node[lab,above] at (-3.4,-5.6) {deploy model};
\end{tikzpicture}
\end{adjustbox}
\caption{Final system architecture: research pipeline and user application.}
\label{fig:arch}
\end{figure}

\subsection{Why the Proposed Design Was Selected}
The proposed design was selected because it balances research reliability,
development effort, scalability, and practicality. Using existing datasets allows
the technical pipeline to be developed early. Creating a smaller custom dataset
makes it possible to focus annotation effort on data specifically relevant to
Bangla propaganda detection. Using an LLM for pre-annotation reduces repetitive
work, while human verification protects the quality of the ground-truth labels.

Similarly, using traditional machine learning as a baseline and transformer models
as the main candidates allows meaningful comparison between simple and advanced
approaches. Starting with text-based detection keeps the initial implementation
manageable, while the architecture leaves room for future multimodal expansion.
The overall design is therefore iterative:
\[\text{Build}\rightarrow\text{Test}\rightarrow\text{Analyze}\rightarrow\text{Improve}\rightarrow\text{Retrain}\rightarrow\text{Re-evaluate}.\]
This makes it possible to improve both the dataset and the model throughout the
project instead of expecting the first version to be perfect.

% ---------------------------------------------------------------------
\section{Project Plan}
\label{sec:plan}
The project lifecycle is planned across three consecutive academic trimesters
(25 weeks), structured into four distinct engineering phases. The schedule is shown
as a Gantt chart in Figure~\ref{fig:gantt}.

\begin{figure}[H]
\centering
\begin{adjustbox}{max width=\textwidth}
\begin{ganttchart}[
    hgrid, vgrid={*{4}{draw=none},dotted},
    x unit=0.42cm, y unit chart=0.6cm, y unit title=0.55cm,
    title height=1, title label font=\tiny\bfseries,
    bar height=0.55, bar label font=\scriptsize,
    group label font=\scriptsize\bfseries,
    bar/.append style={fill=uiublue!70,draw=uiublue},
    group/.append style={fill=uiublue!35,draw=uiublue},
    bar label node/.append style={text width=5.8cm,align=right,anchor=east,font=\scriptsize\RaggedLeft},
    group label node/.append style={text width=5.8cm,align=right,anchor=east,font=\scriptsize\bfseries\RaggedLeft}
  ]{1}{25}
  \gantttitle{Weeks 1--25}{25}\\
  \gantttitlelist{1,...,25}{1}\\
  \ganttgroup{Phase 1: Planning}{1}{8}\\
  \ganttbar{Literature Review \& Taxonomy}{1}{8}\\
  \ganttbar{Requirements \& Architecture Spec.}{3}{8}\\
  \ganttgroup{Phase 2: Data Pipeline}{5}{16}\\
  \ganttbar{Web Scraping \& Data Ingestion}{5}{14}\\
  \ganttbar{LLM Pre-Annotation Prompting}{7}{12}\\
  \ganttbar{Human Expert Annotation \& Verification}{9}{16}\\
  \ganttgroup{Phase 3: Modeling \& Evaluation}{15}{22}\\
  \ganttbar{Classical Baselines (SVM/NB)}{15}{18}\\
  \ganttbar{BanglaBERT \& XLM-R Fine-Tuning}{16}{21}\\
  \ganttbar{Class Imbalance Mitigation \& Testing}{18}{22}\\
  \ganttgroup{Phase 4: UI \& Deployment}{20}{25}\\
  \ganttbar{Explainability Module \& UI Dashboard}{20}{24}\\
  \ganttbar{Benchmarking \& Final Thesis Draft}{22}{25}
\end{ganttchart}
\end{adjustbox}
\caption{Gantt chart of the project execution lifecycle (Weeks 1--25).}
\label{fig:gantt}
\end{figure}

\begin{itemize}
  \item \textbf{Phase 1: Requirements Gathering \& Taxonomy Definition (Weeks
        1--8).} Systematic review of propaganda literature, defining the Bangla
        propaganda technique taxonomy, and architecting the DFD.
  \item \textbf{Phase 2: Data Curation \& Hybrid Annotation (Weeks 5--16).}
        Gathering raw posts, prompt engineering for LLM pre-annotation, and
        conducting human expert verification to construct the verified dataset.
  \item \textbf{Phase 3: Model Training \& Experimental Benchmarking (Weeks
        15--22).} Fine-tuning BanglaBERT and XLM-R, implementing cost-sensitive
        loss weighting, and conducting ablation experiments E1--E6.
  \item \textbf{Phase 4: Dashboard Integration \& Final Thesis (Weeks 20--25).}
        Integrating the prediction engine with the web user interface, generating
        explainability summaries, and writing the final thesis report.
\end{itemize}

\section{Task Allocation}
\label{sec:tasks}
Project responsibilities are divided among team members according to technical
sub-disciplines:
\begin{itemize}
  \item \textbf{Shekh Abdullah Al Mehedi (ID: 112231075):} data pipeline
        architecture, automated web scraping infrastructure, data
        cleaning/normalization, and database management.
  \item \textbf{Mahjabin Khan (ID: 112230177):} annotation pipeline design, LLM
        prompt engineering for pre-annotation, annotator guideline drafting, and
        human verification coordination.
  \item \textbf{Omor Faruck Ullas (ID: 112310384):} machine learning pipeline,
        transformer model fine-tuning (BanglaBERT \& XLM-R), class imbalance
        mitigation, and loss function engineering.
  \item \textbf{Md.\ Rayan Rahman (ID: 112331071):} system evaluation, benchmarking
        experiments (E1--E6), confusion matrix error analysis, and web UI/API
        development.
\end{itemize}

% =====================================================================
%                            CHAPTER 4
% =====================================================================
\chapter{Standards and Design Constraints}
\label{ch:standards}

This chapter details the engineering standards, real-world design constraints,
financial budget analysis, and formal mapping to Washington Accord / BAETE Complex
Engineering Problem Solving criteria guiding the development of the Bangla
propaganda detection framework.

\section{Compliance with Standards}
To ensure reproducibility, maintainability, and interoperability, the project
adheres to established technical standards.

\subsection{Software Standards}
\begin{itemize}
  \item Adherence to \textbf{PEP~8} style guidelines for all Python modules.
  \item Utilization of standardized machine learning frameworks (\textbf{PyTorch},
        \textbf{Hugging Face Transformers}) to ensure cross-platform
        reproducibility.
  \item Implementation of \textbf{RESTful API} conventions using FastAPI/Flask and
        \textbf{JSON Schema} for structured data exchange.
\end{itemize}

\subsection{Hardware and Cloud Standards}
\begin{itemize}
  \item Compliance with standard GPU computing interfaces (CUDA~12.x, OpenCL) to
        allow portable training across both local workstations and cloud
        infrastructure.
\end{itemize}

\subsection{Communication and Security Standards}
\begin{itemize}
  \item Enforcing \textbf{HTTPS (TLS~1.3)} encryption for all web interface
        interactions and API queries.
  \item Adoption of the \textbf{Schema.org ClaimReview} structured data standard to
        ensure interoperability with international fact-checking
        platforms~\cite{google_claimreview}.
\end{itemize}

\section{Realistic Design Constraints}

\subsection{Economic Constraint}
Academic and independent journalistic organizations in Bangladesh operate under
strict financial limits. The framework avoids dependence on proprietary per-token
cloud APIs (e.g., proprietary LLM inference) for runtime detection. Instead, the
final model is fine-tuned into open-weight transformer backbones (BanglaBERT)
capable of low-cost local or edge deployment.

\subsection{Environmental Constraint}
Training large generative architectures from scratch consumes enormous electrical
power. We constrain our carbon footprint by utilizing transfer learning ---
fine-tuning existing open-source pretrained representations rather than training
models from scratch.

\subsection{Ethical Constraint}
Detecting political propaganda carries significant ethical risks regarding
censorship, freedom of speech, and algorithmic political bias. To uphold ethical
standards:
\begin{itemize}
  \item The system outputs probability confidence scores and explainable rationale
        rather than absolute binary judgments.
  \item Training data is balanced across diverse political viewpoints to prevent
        partisan favoritism.
  \item Public data collection strictly complies with platform Terms of Service,
        and personal identifiers are anonymized.
\end{itemize}

\subsection{Health and Safety Constraint}
While software does not cause direct physical harm, coordinated disinformation in
Bangladesh has repeatedly triggered real-world communal violence and riots.
Erroneous detection could amplify panic or silence legitimate reporting. The system
incorporates human-in-the-loop review safeguards before high-consequence moderation
flags are triggered.

\subsection{Social Constraint}
The tool is designed to empower citizens, journalists, and researchers to think
independently rather than having their views manipulated by coordinated actors. The
user interface is crafted in clear, accessible Bangla to serve users across varying
levels of technical literacy.

\subsection{Political Constraint}
The project operates strictly within the legal bounds of Bangladesh's Cyber Security
Act and digital data protection policies, maintaining neutrality and strictly
avoiding algorithmic political profiling.

\subsection{Sustainability}
The framework relies exclusively on open-source libraries and open-weight models,
ensuring that the software can be maintained, adapted, and extended by the research
community without ongoing commercial licensing fees.

\section{Cost Analysis}

\subsection{Proposed Budget}
The proposed budget accounts for dedicated computational resources, data storage,
and verification honorariums (Table~\ref{tab:budget1}).

\begin{table}[H]
\centering
\caption{Proposed budget for the Bangla propaganda detection project.}\label{tab:budget1}
\renewcommand{\arraystretch}{1.3}
\begin{tabularx}{\textwidth}{@{}L{3.6cm}X r@{}}
\toprule
\textbf{Category} & \textbf{Description} & \textbf{Est.\ Cost (BDT)}\\
\midrule
Computational Resources & Dedicated GPU workstation time, RAM, and cloud compute subscriptions for transformer training & 1,80,000\\
Data Collection \& Storage & Cloud database storage for the multimodal social media corpus and automated scraper hosting & 45,000\\
Annotation \& Expert Verification & Honorarium for domain experts and linguists conducting ground-truth label verification & 60,000\\
Miscellaneous \& Contingency & High-speed broadband data packages and unexpected contingency expenses & 5,000\\
\midrule
\multicolumn{2}{@{}l}{\textbf{Total}} & \textbf{2,90,000}\\
\bottomrule
\end{tabularx}
\end{table}

\subsection{Alternate Budget and Rationale}
To accommodate severe financial constraints, an alternate low-cost budget leveraging
free academic resources is established (Table~\ref{tab:budget2}).

\begin{table}[H]
\centering
\caption{Alternate low-cost budget.}\label{tab:budget2}
\renewcommand{\arraystretch}{1.3}
\begin{tabularx}{\textwidth}{@{}L{3.6cm}X r@{}}
\toprule
\textbf{Category} & \textbf{Description} & \textbf{Est.\ Cost (BDT)}\\
\midrule
Computational Resources & Free-tier Google Colab (T4 GPUs) and UIU Department computer-lab workstations & 0\\
Data Collection \& Storage & Local workstation hard-drive storage and university Google Drive allocation & 0\\
Annotation \& Verification & Internal student-team manual verification under faculty supervision & 0\\
Miscellaneous & Contingency buffer & 2,500\\
\midrule
\multicolumn{2}{@{}l}{\textbf{Total}} & \textbf{2,500}\\
\bottomrule
\end{tabularx}
\end{table}

\textbf{Rationale.} While the alternate budget incurs near-zero financial cost, the
Proposed Budget (Table~\ref{tab:budget1}) is strongly preferred. Training deep
transformer networks on free Colab instances suffers from frequent kernel
disconnects, timeouts, and limited GPU memory. Dedicated resources guarantee
reliable experimental reproducibility.

\subsection{Sustainability and Post-FYDP Revenue Model}
To sustain the platform after project completion:
\begin{enumerate}
  \item \textbf{Open-Source Core:} the trained model weights and benchmark dataset
        are freely released to the academic community.
  \item \textbf{B2B / API Licensing for Newsrooms:} offering an automated
        verification API to independent digital newsrooms and fact-checking
        organizations in Bangladesh under a software-as-a-service (SaaS)
        subscription model.
  \item \textbf{Grant Funding:} seeking research grants from organizations focused
        on digital literacy, counter-disinformation, and low-resource NLP.
\end{enumerate}

\section{Complex Engineering Problem (BAETE / Washington Accord)}

\subsection{Complex Problem Solving Mapping}
The project is mapped against the Washington Accord complex engineering problem
criteria (P1--P7) in Table~\ref{tab:p1p7}.

\begin{table}[H]
\centering
\caption{Mapping with complex problem solving criteria (P1--P7).}\label{tab:p1p7}
\renewcommand{\arraystretch}{1.5}
\resizebox{\textwidth}{!}{%
\begin{tabular}{@{}c|c|c|c|c|c|c|c@{}}
\toprule
\textbf{Attribute} & \makecell{\textbf{P1}\\Depth of\\Knowledge} & \makecell{\textbf{P2}\\Conflicting\\Requirements} & \makecell{\textbf{P3}\\Depth of\\Analysis} & \makecell{\textbf{P4}\\Familiarity\\of Issues} & \makecell{\textbf{P5}\\Applicable\\Codes} & \makecell{\textbf{P6}\\Stakeholder\\Involvement} & \makecell{\textbf{P7}\\Inter-\\dependence}\\
\midrule
\textbf{Status} & \ck & \ck & \ck & \ck & \ck & \ck & \ck\\
\bottomrule
\end{tabular}}
\end{table}

\begin{itemize}
  \item \textbf{P1: Depth of Knowledge Required.} Synthesizes Natural Language
        Processing (transformers, attention mechanisms), multimodal computer vision
        (ViT, ResNet), graph theory (coordinated social networks), and linguistic
        rhetoric. Aligns with \textbf{WK3} (advanced engineering knowledge) and
        \textbf{WK4} (specialist AI/ML architecture knowledge).
  \item \textbf{P2: Range of Conflicting Requirements.} Balancing deep contextual
        accuracy with computational inference latency; balancing aggressive
        propaganda detection with the protection of free speech and political
        satire.
  \item \textbf{P3: Depth of Analysis Required.} In-depth modeling of subtle
        rhetorical propaganda techniques (loaded language, fear appeals,
        bandwagon) in morphologically rich, code-mixed Bangla text.
  \item \textbf{P4: Familiarity of Issues.} Detecting coordinated propaganda
        campaigns in under-resourced Bangla digital media is largely uncharted
        territory compared to standard English fake-news benchmarks.
  \item \textbf{P5: Extent of Applicable Codes.} Strict adherence to data privacy,
        digital copyright, platform terms of service, and ethical AI fairness
        standards.
  \item \textbf{P6: Extent of Stakeholder Involvement.} Involves diverse
        stakeholders including digital media consumers, journalists, academic
        researchers, and content moderators.
  \item \textbf{P7: Inter-dependence.} High inter-dependence across web scraping,
        preprocessing, hybrid annotation, transformer fine-tuning, and user result
        generation; an error or bias in annotation ripples through the entire
        classification pipeline.
\end{itemize}

\subsection{Complex Engineering Activities Mapping}
The project is mapped against the complex engineering activity attributes
(A1--A5) in Table~\ref{tab:a1a5}.

\begin{table}[H]
\centering
\caption{Mapping with complex engineering activities (A1--A5).}\label{tab:a1a5}
\renewcommand{\arraystretch}{1.5}
\begin{tabular}{@{}c|c|c|c|c|c@{}}
\toprule
\textbf{Attribute} & \makecell{\textbf{A1}\\Range of\\Resources} & \makecell{\textbf{A2}\\Level of\\Interaction} & \makecell{\textbf{A3}\\Innovation} & \makecell{\textbf{A4}\\Consequences\\for Society} & \makecell{\textbf{A5}\\Familiarity}\\
\midrule
\textbf{Status} & \ck & \ck & \ck & \ck & \ck\\
\bottomrule
\end{tabular}
\end{table}

\begin{itemize}
  \item \textbf{A1: Range of Resources.} GPU clusters, cloud storage, large-scale
        web scrapers, transformer toolkits, and human expert linguistic annotators.
  \item \textbf{A2: Level of Interaction.} Continuous synergy between computational
        NLP subsystems and human domain experts during the two-tier verification
        workflow.
  \item \textbf{A3: Innovation.} Novel hybrid LLM--human annotation paradigm and
        multi-faceted coordinated campaign detection framework tailored
        specifically to Bangla online media.
  \item \textbf{A4: Consequences for Society and Environment.} Substantial societal
        impact in safeguarding civic discourse and countering dangerous communal
        manipulation in Bangladesh.
  \item \textbf{A5: Familiarity.} Requires innovative adaptation of cross-lingual
        transfer learning to overcome the severe data scarcity characteristic of
        low-resource Bangla NLP.
\end{itemize}

\section{Summary}
This chapter formalized the engineering standards, realistic constraints, economic
feasibility budgets, and rigorous BAETE accreditation mappings that validate this
research as a comprehensive complex engineering capstone project.

% =====================================================================
%                            BIBLIOGRAPHY
% =====================================================================
\clearpage
\phantomsection
\addcontentsline{toc}{chapter}{Bibliography}
\begin{thebibliography}{99}
\setlength{\itemsep}{3pt}

\bibitem{hristakieva2022} K.~Hristakieva, S.~Cresci, G.~Da San Martino, M.~Conti, and P.~Nakov, ``The spread of propaganda by coordinated communities on social media,'' in \emph{Proceedings of the 14th ACM Web Science Conference 2022 (WebSci '22)}. Association for Computing Machinery, 2022, pp.~191--201.

\bibitem{hossain2020} M.~Z. Hossain, M.~A. Rahman, M.~S. Islam, and S.~Kar, ``BanFakeNews: A dataset for detecting fake news in Bangla,'' in \emph{Proceedings of the Twelfth Language Resources and Evaluation Conference}. European Language Resources Association, 2020, pp.~2862--2871.

\bibitem{kamruzzaman2023} M.~Kamruzzaman, M.~M.~I. Shovon, and G.~Kim, ``BanMANI: A dataset to identify manipulated social media news in Bangla,'' in \emph{Proceedings of the Workshop on Computational Terminology in NLP and Translation Studies (ConTeNTS) Incorporating the 16th Workshop on Building and Using Comparable Corpora (BUCC)}. Shoumen, Bulgaria: INCOMA Ltd., 2023, pp.~51--58.

\bibitem{ng2023} K.~W. Ng and A.~Iamnitchi, ``Coordinated information campaigns on social media: A multifaceted framework for detection and analysis,'' in \emph{Disinformation in Open Online Media}. Springer Nature Switzerland, 2023, pp.~103--118.

\bibitem{duridi2025} T.~Duridi, L.~Atwe, A.~Jaber, E.~Daraghmi, and P.~Mart\'inez, ``Detection of propaganda and bias in social media: A case study of the Israel--Gaza war (2023),'' in \emph{Proceedings of the 2025 International Conference on New Trends in Computing Sciences (ICTCS)}, 2025.

\bibitem{plikynas2025} D.~Plikynas, I.~Rizgelien\.e, and G.~Korvel, ``Systematic review of fake news, propaganda, and disinformation: Examining authors, content, and social impact through machine learning,'' \emph{IEEE Access}, vol.~13, pp.~17\,583--17\,629, 2025.

\bibitem{lora2023} S.~K. Lora, I.~Jahan, R.~Hussain, R.~Shahriyar, and A.~B.~M.~A. Al Islam, ``A transformer-based generative adversarial learning to detect sarcasm from Bengali text with correct classification of confusing text,'' \emph{Heliyon}, vol.~9, no.~12, 2023.

\bibitem{hasan2026} M.~A. Hasan, F.~Alam, M.~F. Hossain, U.~Naseem, and S.~I. Ahmed, ``LLM-based multi-task Bangla hate speech detection: Type, severity, and target,'' in \emph{Proceedings of the 64th Annual Meeting of the Association for Computational Linguistics (ACL)}, 2026, pp.~33\,962--33\,980.

\bibitem{ahmad2025} P.~N. Ahmad, J.~Guo, N.~M. AboElenein, Q.~M. ul~Haq, S.~Ahmad, A.~D. Algarni, and A.~A. Ateya, ``Hierarchical graph-based integration network for propaganda detection in textual news articles on social media,'' \emph{Scientific Reports}, vol.~15, 2025.

\bibitem{lin2024} S.-Y. Lin, Y.-C. Chen, Y.-H. Chang, S.-H. Lo, and K.-M. Chao, ``Text-image multimodal fusion model for enhanced fake news detection,'' \emph{Science Progress}, 2024.

\bibitem{hee2022} M.~S. Hee, R.~K.-W. Lee, and W.-H. Chong, ``On explaining multimodal hateful meme detection models,'' in \emph{Proceedings of the ACM Web Conference (WWW)}, 2022, pp.~3651--3655.

\bibitem{declerck2024} B.~De~Clerck, J.~C. Fernandez~Toledano, F.~Van Utterbeeck, and L.~E.~C. Rocha, ``Detecting coordinated and bot-like behavior in Twitter: The J\"urgen Conings case,'' \emph{EPJ Data Science}, vol.~13, no.~1, 2024.

\bibitem{sela2025} A.~Sela, O.~Neter, V.~Lohr, P.~Cihelka, F.~Wang, M.~Zwilling, J.~P. Sabou, and M.~Ulman, ``Signals of propaganda---detecting and estimating political influences in information spread in social networks,'' \emph{PLOS ONE}, vol.~20, 2025.

\bibitem{schneider2026} P.~J. Schneider, L.~Yuan, and M.-A. Rizoiu, ``Beyond content: Behavioral policies reveal actors in information operations,'' \emph{npj Complexity}, 2026.

\bibitem{romim2022} N.~Romim, M.~Ahmed, M.~S. Islam, A.~Sen~Sharma, H.~Talukder, and M.~R. Amin, ``BD-SHS: A benchmark dataset for learning to detect online Bangla hate speech in different social contexts,'' in \emph{Proceedings of the 13th Language Resources and Evaluation Conference (LREC)}, 2022, pp.~5153--5162.

\bibitem{google_factcheck} Google, ``Fact check tools API,'' Google for Developers. [Online]. Available: \url{https://developers.google.com/fact-check/tools/api}

\bibitem{google_claimreview} Google, ``Fact check (ClaimReview) structured data,'' Google Search Central. [Online]. Available: \url{https://developers.google.com/search/docs/appearance/structured-data/factcheck}

\bibitem{haider2020} S.~Haider, L.~Luceri, A.~Deb, A.~Badawy, N.~Peng, and E.~Ferrara, ``Detecting social media manipulation in low-resource languages,'' 2020.

\bibitem{wang2022} M.-H. Wang, W.-Y. Chang, K.-H. Kuo, and K.-Y. Tsai, ``Analyzing image-based political propaganda in referendum campaigns: From elements to strategies,'' 2022.

\bibitem{pina2025} C.~A. Pi\~na-Garc\'ia, ``In-context learning for propaganda detection on Twitter Mexico using large language model Meta AI,'' \emph{Telematics and Informatics Reports}, vol.~19, 2025.

\bibitem{gaeta2025} A.~Gaeta, V.~Loia, A.~Lorusso, F.~Orciuoli, and A.~Pascuzzo, ``Towards a LLM-based intelligent system for detecting propaganda within textual content,'' \emph{Computers \& Electrical Engineering}, vol.~128, 2025.

\bibitem{maarouf2024} A.~Maarouf, D.~B\"ar, D.~Geissler, and S.~Feuerriegel, ``HQP: A human-annotated dataset for detecting online propaganda,'' in \emph{Findings of the Association for Computational Linguistics: ACL 2024}, 2024, pp.~6064--6089.

\end{thebibliography}

\end{document}